\chapter{Introduction}

\section{Parthenium as an invasive alien plant species}

Invasive alien plant species (IAPS) pose severe threats to South Asian biodiversity, agriculture, and human health, with \textit{Parthenium hysterophorus} among the most aggressive invaders [X]. Native to the Americas, this annual herb, commonly known as ``Congress grass'' or ``Gajar Booti'', has rapidly colonised diverse habitats across Asia, Africa, and Australia [X]. Its success stems from high phenotypic plasticity, allelopathic suppression, and broad tolerance, together with a persistent soil seed bank capable of surviving wide thermal ranges [X]. Despite decades of spread since its mid-twentieth-century introduction to South Asia [X], key gaps remain regarding how microclimate, human dispersal, and land-use changes interact to drive its expansion. This study investigates the spatiotemporal spread of \textit{P. hysterophorus} in India and Pakistan, mapping environmental and anthropogenic drivers to address management gaps and inform biosecurity policy [X].

Parthenium exemplifies a fast-maturing annual plant characterised by a deep taproot and an erect stem. Most individuals remain below 1.5 m in height, although some may reach up to 2.5 m [X]. The small, white flowers (approximately 4 mm in diameter) display five distinct lobes, and each plant is capable of producing up to 20{,}000 viable seeds from its 2 mm wedge-shaped achenes [X]. The soil seed bank may contain over 340 million seeds per hectare. Rapid proliferation during germination, combined with fast growth and allelopathic properties, enables \textit{Parthenium} to establish strongly, suppress neighbouring vegetation, and maintain a persistent seed bank [X].

\section{Introducing the spread of parthenium weed in South Asia}

\textit{Parthenium hysterophorus} was likely introduced into India in the 1950s through contaminated PL 480 wheat shipments from the United States, initially appearing as stray plants in Pune [X]. Although records from 1814 suggest earlier localised occurrences, large-scale expansion began in the mid-twentieth century [X]. By 1972, the weed had spread from Kashmir to Kerala, and by 1979 it had reached Assam. It now occurs in every Indian state, with Karnataka alone accounting for approximately 5 million hectares of infestation [X]. In Pakistan, \textit{Parthenium} was first confirmed in the Gujrat district of Punjab in 1980, likely entering through cross-border trade at the Wagah border [X]. It subsequently spread across Punjab, Khyber Pakhtunkhwa (KP), Islamabad Capital Territory (ICT), and Azad Jammu and Kashmir, with severe infestations concentrated in central and northern Punjab, ICT, and KP, and the invasion continuing southward [X]. \textit{Parthenium} now occurs throughout South Asia and is among the region's most problematic weeds.

This dissertation examines the environmental, climatic, and anthropogenic factors underlying the transition of \textit{P. hysterophorus} from localised introductions to widespread subcontinental invasion. Understanding these drivers is essential for explaining its ecological impacts and the persistent challenges of containment across South Asia.

\textit{Parthenium} spreads through both natural and human-mediated pathways [X]. Farmers identify water and vehicles as major dispersal agents, while wind, animals, and contaminated agricultural products facilitate movement of its lightweight seeds [X]. Human-mediated transport enables the species to overcome barriers such as rivers, mountains, and climatic zones. This study therefore investigates the dispersal dynamics and transport networks, including road hierarchies, that facilitate its regional spread.

Despite extensive management efforts, including mechanical removal, chemical herbicides, competitive replacement, and biological control, \textit{Parthenium} remains highly invasive. Prolific seed production, a persistent soil seed bank, rapid growth, broad ecological tolerance, and the ability to exploit disturbed habitats promote rapid establishment. Through allelopathy and competition for resources, it suppresses native vegetation, reducing biodiversity and promoting dense monocultures. These traits are particularly advantageous across South Asia's agricultural land, grazing areas, roadsides, and other disturbed habitats.

The region's warm climate, seasonal rainfall, extensive agriculture, and habitat disturbance further facilitate establishment and persistence. In India and Pakistan, cultivated landscapes, irrigation, livestock movement, transport corridors, and degraded habitats form interconnected pathways for dispersal. Together, these environmental and anthropogenic factors have enabled the rapid expansion of \textit{P. hysterophorus} across South Asia, highlighting the need to understand the mechanisms sustaining its invasion.

\section{Harmful effects of \textit{Parthenium hysterophorus}}

\textit{Parthenium hysterophorus} rapidly colonises disturbed habitats, including wastelands, roadsides, waterways, and agricultural fields, across a wide range of environmental conditions. Prolific seed production, rapid growth, and broad climatic tolerance enable it to form dense stands that displace native species and disrupt ecosystem processes. Disturbance, rainfall, and temperature influence its distribution and impacts in South Asia, while land-use change and soil disturbance further enhance persistence and agricultural impacts through allelopathy and competition for resources.

Human activity, particularly agriculture and transportation, also drives its spread and exposure. Rural communities and agricultural workers are disproportionately affected, as infestations frequently occur around cultivated land and transport networks. Exposure to the plant's allergenic compounds can cause dermatitis and respiratory conditions, including bronchitis. Increasing land use and connectivity therefore amplify risks to both agricultural production and public health [X].

The impacts of \textit{P. hysterophorus} extend across interconnected ecological, agricultural, and socioeconomic systems. These effects are particularly pronounced in South Asia, where high population density, agricultural dependence, and favourable environmental conditions facilitate rapid spread and exposure. Consequently, \textit{P. hysterophorus} represents not only an ecological invasion but also a significant agricultural, economic, and public health concern, reinforcing the need for integrated regional management.

\section{Biocontrol management and climate expansion}

Long-term management of \textit{Parthenium hysterophorus} remains challenging due to its adaptability and ongoing expansion. Rising atmospheric CO$_2$ may further increase herbicide resistance, necessitating a reassessment of traditional control strategies. Climate change is projected to facilitate additional spread of \textit{Parthenium} and other invasive species, amplifying the urgency for proactive management.

Species distribution models reinforce these risks: CLIMEX analyses indicate much of South Asia is already suitable for \textit{Parthenium} establishment. Climate change scenarios predict continued northward expansion, with irrigation and warming enabling colonisation of previously unsuitable regions. Identifying high-risk areas will be essential for targeted monitoring and early intervention. Control methods, including burning, herbicides, and biological agents such as beetles, moths, weevils, and fungi, have yielded variable results. The biocontrol beetle \textit{Zygogramma bicolorata}, for example, is highly sensitive to temperatures above 35$^{\circ}$C, complicating sustainable management as climate warms [X].

\section{Study aims}

Despite extensive research, major gaps remain, particularly concerning cross-border dynamics and climate-driven changes in distribution. India and Pakistan are considered together here because of their shared invasion history, contiguous borders, and differing climates and agricultural systems that distinctly influence \textit{Parthenium}'s distribution. However, comparative region-wide analyses are scarce, limiting understanding of how climate change may affect both \textit{Parthenium} and its biocontrol beetle across these countries.

This study models current and projected distributions of \textit{Parthenium} and its biocontrol beetle across India and Pakistan, integrating occurrence records with environmental variables. It hypothesises that climate change will alter both distributions, with direct management implications. By identifying emerging invasion risks and changes in habitat suitability, this research highlights where management, especially biological control, must adapt, guiding proactive strategies for vulnerable regions.